% Please do not change %
\documentclass[11pt]{article} % font size
\usepackage[margin=1in]{geometry} % margins
\usepackage{amsmath,amsthm,amssymb} % for the  common "math symbols"
\usepackage{algorithm}
\usepackage{algpseudocode}
\usepackage[shortlabels]{enumitem} % for enumeration
\usepackage{booktabs} % if you need tables
\usepackage{listings}
\usepackage{color}
\DeclareMathOperator*{\argmax}{argmax}
\DeclareMathOperator*{\argmin}{argmin}

\lstset{
  frame=none,
  xleftmargin=2pt,
  stepnumber=1,
  numbers=left,
  numbersep=5pt,
  numberstyle=\ttfamily\tiny\color[gray]{0.3},
  belowcaptionskip=\bigskipamount,
  captionpos=b,
  escapeinside={*'}{'*},
  language=haskell,
  tabsize=2,
  emphstyle={\bf},
  commentstyle=\it,
  stringstyle=\mdseries\rmfamily,
  showspaces=false,
  keywordstyle=\bfseries\rmfamily,
  columns=flexible,
  basicstyle=\small\sffamily,
  showstringspaces=false,
  morecomment=[l]\%,
}
% Feel free to include additional packages (if needed), but make sure it does not override the margin/font requirements.
%
\setlength{\parindent}{0pt} % no indent for new paragraphs.
\setlength{\parskip}{5pt} % distance between paragraphs, which will be used when you have a blank line in your LaTeX code.


%%%%%%%%%%%%%%%%%%%%%%%%%%%%%%%%%%%%%%%%%%%%%%%%%%%%%%%%%%%%%%%%%%%%%%%%%%%%%%
\title{Probability and Computing: \\Exercise Sheet 6}
%
\author{Wira Azmoon Ahmad}
\date{24 Nov 2023}
%
\begin{document}
%
\maketitle

%You may use the usual \LaTeX\ environments to structure your answers:
%\begin{align}
%        \label{basegnn}
%        \mathbf{h}_u^{(t+1)} = \sigma \bigg( \mathbf{W}_\text{self}^{(t)} \mathbf{h}_u^{(t)}  + \mathbf{W}_\text{neigh}^{(t)} \sum_{v \in N(u)} \mathbf{h}_v^{(t)} + \mathbf{b}^{(t)} \bigg). 
%\end{align}


%\begin{table}[h]
%\caption{A simple table.}
%\centering
%\begin{tabular}[t]{lrrrr}
%\toprule
%Graphs & Can & Be  & Fun & Too \\ 
%\midrule 
%$G_1$ & $1$  & $2$ & $3$ & $4$ \\ 
%$G_2$ & $-1$  & $-2$ & $-3$ & $-4$ \\ 
%$G_3$ & $0.1$  & $0.2$ & $0.3$ & $0.4$ \\ 
%\bottomrule
%\end{tabular}
%\label{tab}
%\end{table}

\section*{Question 1}

% To answer each part of the question, please simply use a paragraph command 
% \paragraph{(a)}
Let $(X_t,Y_t)$ be a coupling of the Markov chain $M_t$.
We choose a coupling as follows.

\begin{itemize}
    \item Choose positions $j_1,j_2$ u.a.r. from $\{1,...,n\}$.
    \item Obtain $X_{t+1}$ from $X_t$ by swapping the card at position $j_1$ with card at position $j_2$.
    Call the card at position $j_2$ card $C$.
    \item Obtain $Y_{t+1}$ from $Y_t$ by swapping the card at position $j_1$ with card $C$.
\end{itemize}

We show that $Y_t$ is a faithful copy of the original chain.
Consider any position $k_1,k_2$. Let $C'$ be the card in position $k_2$ in $Y_t$.
Let $j_2$ be the position of $C'$ in $X_t$, and let $j_1=k_1$.
The probability that $j_1,j_2$ are chosen are each $1/n$, 
so this is the probability that position $k_1,k_2$ are chosen in the transition from $Y_t$ to $Y_{t+1}$.

% Let $d(x,y)$ be the number of positions on which the states $x,y$ disagree.
% Suppose $d(x,y)=k$. At each time step, if $C_i$ in the same position for both, 
% or if the card at $j_i$ is the same for both $X_t$ and $Y_t$,
% then the number of disagreements stay the same.
% Otherwise, it goes down by at least one, with probability $(d(x,y)/n)^2$ since the events are independent.
% Then for $d(x,y)\geq 1$,
% \[\mathbb{E}[d(X_1,Y_1)|(X_0,Y_0)=(x,y)] \leq d(x,y) - \left(\frac{d(x,y)}{n}\right)^2 
% \leq d(x,y) - d(x,y)/n^2 \leq\beta d(x,y)\]
% for $\beta = 1-1/n^2$.


Let $Z_t$ denote the random variable of the number of cards whose positions are different for $X_t$ and $Y_t$.
Clearly $0\leq Z_t \leq n$. At each time step $t$, we choose a (card, position) pair $(C_i,j_i)$.
If $C_i$ in the same position for both, or if the card at $j_i$ is the same 
for both $X_t$ and $Y_t$, then $Z_{t+1}=Z_t$. 
Otherwise, we have $Z_{t+1} \leq Z_t -1$.
The probability of the chosen card $C_i$ being in different positions is $Z_t/n$,
and the probability of the chosen position $j_i$ having different cards is also $Z_t/n$.
Since both events are independent, 
\[\Pr(Z_{t+1} \leq Z_t - 1) = \left(\frac{Z_t}{n}\right)^2\]

Let $T_i$ be the time spent for which $Z_t=i$. 
Then the maximum total time spent before coupling is $T\leq\sum_{i=1}^n T_i$, 
then $\mathbb{E}[T]\leq\sum_{i=1}^n \mathbb{E}[T_i]$.
$T_i$ is a geometric r.v. with parameter $p_i=\frac{i^2}{n^2}$, so $\mathbb{E}[T_i] = \frac{n^2}{i^2}$.
Then, the expected time spent  is
\[\mathbb{E}[T]\leq n^2\sum_{i=1}^{n} \frac{1}{i^2}< \frac{n^2\pi^2}{6}\]

By Markov's inequality, the probability not to obtain the desired sequence after $\frac{n^2\pi^2}{3}$
steps is at most $1/2$.
Thus, after $\frac{kn^2\pi^2}{3}$ steps, the probability not to obtain the sequence is at most $2^{-k}$.
We deduce that $\Pr(X_t\neq Y_t)\leq \varepsilon$ if $t\geq \frac{\log_2(1/\varepsilon)n^2\pi^2}{3}$,
so by the Coupling Lemma, the mixing time is
\[\tau(\varepsilon) = O(n^2 \log(1/\varepsilon))\]

\section*{Question 2}
The Markov chain is \textbf{finite}, with state space $\Omega=\mathcal{I}(G)\subseteq 2^V$.
It is \textbf{irreducible}, since we can get from any independent sets 
$I\in\mathcal{I}(G)$ to $I'\in\mathcal{I}(G)$ by adding or removing vertices approriately.
It is \textbf{aperiodic}, since it can self-loop. 
Thus, there is a unique stationary distribution $\pi$.

Let $n=|V|$. Observe that for $I,J\in\mathcal{I}(G),I\neq J$,
\begin{itemize}
    \item $P_{I,J} = \frac{\lambda}{2n(1+\lambda)}$, if $J\in\mathcal{I}(G)$ has 1 more vertex than $I$
    \item $P_{I,J} = \frac{1}{2n(1+\lambda)}$, if $J\in\mathcal{I}(G)$ has 1 less vertex than $I$
    \item $P_{I,J} = 0$, otherwise
\end{itemize}
Consider the distribution of the Gibbs measure, so $\pi_I = \lambda^{|I|}/Z(G)$.
Consider also arbitrary $I,J\in\mathcal{I}(G)$, where $|I|+1=|J|$.

\begin{equation}
\begin{aligned}
\pi_I P_{I,J} &= \frac{\lambda^{|I|}}{Z(G)}\cdot\frac{\lambda}{2n(1+\lambda)}\\
&=\frac{\lambda^{|I|+1}}{Z(G)}\cdot\frac{1}{2n(1+\lambda)}\\
&=\frac{\lambda^{|J|}}{Z(G)}\cdot\frac{1}{2n(1+\lambda)}\\
&=\pi_J P_{J,I}
\end{aligned}
\end{equation}
Clearly if $P_{I,J} = P_{J,I} = 0, \pi_I P_{I,J} = 0 = \pi_J P_{J,I}$.
Thus, $\forall I,J, \pi_I P_{I,J} = \pi_J P_{J,I}$.
By time-reversibility in Theorem 4.2.12, the Gibbs measure is the stationary distribution $\pi$.

Suppose the state space $\Omega$ is instead the set of all subgraphs of $G$.
The chain still has the same unique distribution $\pi$ being the Gibbs measure on the independent sets.
If a state $I\notin \mathcal{I}(G)$, then the algorithm will have the state only remain the same
or remove vertices until it is an independent set, after which, it will not return to being a non-independent set.

Now let $H$ be the graph on $\Omega$ in which there is an edge from $x$ to $y$ if and only if
$x$ and $y$ differ by only 1 vertex. For all edges $(x,y), d(x, y)=1$.
Since $\Omega$ contains all subgraphs, the corresponding distance metric $d$ 
is just number of vertices that are in either $X_t$ or $Y_t$, and not both.
For all states $(x,y), d (x,y)$ is just $|x-y|+|y-x|$. 
We choose the coupling where both $(X_t), (Y_t)$ make the same decision in the original Markov chain.

\clearpage
Take $t=0$ and suppose that $(X_0,Y_0)=(x,y)$ where $(x,y)$ is an edge of $H$ so $x$ and $y$
differ by only 1 vertex.
WLOG, suppose $x=I,y=I\cup \{u\}$. Let $d_t=d(X_t,Y_t)$, and let $d_0=1$.
If $X_1 \leftarrow X_0$, so $Y_1 \leftarrow Y_0$, then $d_1=d_0=1$.
Otherwise, consider when the vertex $v$ is chosen u.a.r. from $V$.

\begin{itemize}
    \item If $v\in X_0$, so $v\in Y_0$, then $d_1-d_0=0$ unchanged for any decision made.
    \item Else If $v=u$, then $d_1-d_0=-1$ decrease for any decision made.
    \item Else If $v\in N(u)$,
    \begin{itemize}
        \item Removing $v$ will leave $d_1-d_0=0$ unchanged.
        \item Adding $v$ will have $d_1-d_0=1$ increase, since it can only be added to $X_0$, and not $Y_0$.
    \end{itemize}
    \item Else If $v \notin N(u)$, then $d_1-d_0=0$ unchanged for any decision made.
\end{itemize}

Analysing all cases, only 2 of them will cause $d_1 \neq d_0$, and noting $d_0=1$, so
\begin{equation}
\begin{aligned}
\mathbb{E}[d_1] &=1-\frac{1}{2n}+\frac{\lambda}{1+\lambda}\cdot \frac{d(u)}{2n}\\
&\leq 1 -\frac{1}{2n}+\frac{\lambda}{1+\lambda}\cdot \frac{3}{2n}\\
&=1 + \frac{1}{2n}\left(-1+\frac{3\lambda}{1+\lambda}\right)\\
&=1 + \frac{1}{2n}\cdot\frac{2\lambda - 1}{1+\lambda}\\
&=d_0\left(1 + \frac{1}{2n}\cdot\frac{2\lambda - 1}{1+\lambda}\right)\\
&< d_0
\end{aligned}
\end{equation}
Note that since $0<\lambda<\frac{1}{2}, -1<\frac{2\lambda - 1}{1+\lambda}<0$.

Now we can apply the path coupling lemma with 
$\beta = \left(1 + \frac{1}{2n}\cdot\frac{2\lambda - 1}{1+\lambda}\right)$
and $D=n$, since the maximum the maximum number of vertices $X_t,Y_t$ can differ by is $n$.
This shows that the mixing time satisfies $\tau(\varepsilon)=O(\ln(n/\varepsilon)/\ln(1/\beta))$.
Since $1+x\geq e^x$, and $1+\lambda=\Theta(1)$,
\[\beta \leq e^{\frac{1}{2n}\cdot\frac{2\lambda-1}{1+\lambda}}\]
\[\frac{1}{\beta} \geq e^{-\frac{1}{2n}\cdot\frac{2\lambda-1}{1+\lambda}}\]
\[\ln\frac{1}{\beta} \geq -\frac{1}{2n}\cdot\frac{2\lambda-1}{1+\lambda}\]
\[\frac{\ln(\frac{n}{\varepsilon})}{\ln\frac{1}{\beta}} 
\leq \frac{\ln(\frac{n}{\varepsilon})}{\frac{1-2\lambda}{2n(1+\lambda)}}
= \frac{2n(1+\lambda)\ln(\frac{n}{\varepsilon})}{1-2\lambda}\]
\[\tau(\varepsilon)=O\left(\frac{n\ln(\frac{n}{\varepsilon})}{1-2\lambda}\right)\]

\section*{Question 3}
Let $X_t$ denote the random variable where $X_t=1$ and $X_t=-1$ 
if the clockwise and anticlockwise are taken at step $t\geq 1$, and $X_0=0$.
Then $Z_t=\sum_{i=0}^t X_t$ is the position of the fox at step $t$.

% Let $X_t$ denote the position of the wolf at time $t$.
% Let $Z_t$ denote the martingale where $Z_{t+1}-Z_t=1$ if step $t$ is clockwise, and $Z_{t+1}-Z_t=-1$ if step $t$ is an
We treat going anticlockwise from 0 as going into the negative values to the minimum of $-(n-1)$.
Clearly we have $|Z_t|\leq n-1$.
Then the event that the fox eats sheep $i$ last is the union of the events 
that it reaches the sheep going clockwise, or anticlockwise, which are clearly

Consider the fox eats the sheep clockwise.
Then it must have reached sheep $i+1$ first before reaching $i$,
followed by reaching $i$ from $i-1$ before reaching $i$ from $i+1$.
Then this is the same reaching $i+1-n$ before $i$ in random walk on a line in $[i+1-n,i]$, 
followed by reaching $n-1$ before $-1$ in a random walk on a line in $[-1,n-1]$

Similarly, if the fox eats the sheep anticlockwise,
it reached the sheep $i-1$ first before reaching $i$,
followed by reaching $i$ from $i+1$ before reaching $i$ from $i-1$.
This is the same as reaching $i-1$ before $i-n$ in random walk on a line in $[i-n, i-1]$,
followed by reaching $1-n$ before $1$ in a random walk on a line in $[1-n,1]$.

Both these events are then consecutive instances of the gambler's fortune,
with $a_{1,1}=i,b_{1,1}=n-i-1,a_{1,2}=n-1,b_{1,2}=1$ in the clockwise case,
and $a_{2,1}=i-1,b_{2,1}=n-i,a_{2,2}=1,b_{2,2}=n-1$. 
Note that for all cases $a_{i,1}+b_{i,1}=n-1$, and $a_{i,2}+b_{i,2}=n$.

Since we the know the probability of reaching $a$ before $b$ is $\frac{b}{a+b}$,
and $b$ before $a$ is $\frac{a}{a+b}$,
the probability the fox eats sheep $i$ last is
\[\frac{a_{1,1} b_{1,2}}{n(n-1)} + \frac{b_{2,1} a_{2,2}}{n(n-1)}
=\frac{i+n-i}{n(n-1)}
=\frac{1}{n-1}
\]
This is independent of $i$, so every sheep is equally likely to be eaten last.

\clearpage
\section*{Question 4}
Note that
\[Z_j - Z_{j-1} = \sum_{i=1}^j X_i - \sum_{i=1}^{j-1} X_i = X_j\]
Then
\[Y_i = \sum_{j=1}^i \sum_{k=1}^{j-1} X_k X_j\\
= \sum_{j=1}^i X_j \sum_{k=1}^{j-1} X_k\]
so $Y_n$ is a function of $X_0,...,X_n$, and
\begin{equation}
\begin{aligned}
\mathbb{E}[|Y_n|]&=\mathbb{E}\left[\left|\sum_{i=1}^n X_i Z_{i-1}\right|\right]\\
&=\mathbb{E}\left[\left|\sum_{i=1}^n Z_{i-1}\right|\right] 
\quad\text {(since $X_i\in\{-1,1\}$)}\\
&\leq \sum_{i=1}^n \mathbb{E}[|Z_{i-1}|]\\
&\leq \sum_{i=1}^n (i-1)\\
&<\infty
\end{aligned}
\end{equation}

We prove by induction that 
\[ \mathbb{E}[Y_{n+1}|X_1,...,X_n]=Y_n\]
Base case: $n=0$
\begin{equation}
\begin{aligned}
\mathbb{E}[Y_2|X_1]&=\mathbb{E}[X_1 Z_0 + X_2 Z_1 | X_1]\\
&= \mathbb{E}[X_1 Z_0 | X_1] + \mathbb{E}[X_2 Z_1 | X_1] \quad\text{(linearity of expectation)}\\
&= 0 + \mathbb{E}[X_2 X_1 | X_1] \quad\text{(since $Z_0=0$)}\\
&= \mathbb{E}[X_2]X_1 \quad\text{(since $X_1,X_2$ independent)}\\
&= 0 \quad\text{(since $\mathbb{E}[X_i]=0$)}\\
&= Y_1
\end{aligned}
\end{equation}

Suppose
\[\mathbb{E}[Y_n|X_1,...,X_{n-1}]=Y_{n-1}\]

Note that
\[Y_{n+1} = Y_n + X_{n+1} Z_n = Y_{n-1} + X_n Z_{n-1} + X_{n+1} Z_n\]
\begin{equation}
\begin{aligned}
\mathbb{E}[Y_{n+1}|X_1,...,X_n] &= \mathbb{E}[Y_{n-1} + X_n Z_{n-1} + X_{n+1}Z_n|X_1,...,X_n]\\
&= \mathbb{E}[Y_{n-1}|X_1,...,X_n] + \mathbb{E}[ X_n Z_{n-1}|X_1,...,X_n] + \mathbb{E}[X_{n+1}Z_n|X_1,...,X_n]\\
&= \mathbb{E}[Y_{n-1}|X_1,...,X_{n-1}] + X_n Z_{n-1} + Z_n\mathbb{E}[X_{n+1}]\\
&\quad\quad\text{(since $Y_{n-1}$ ind. $X_n$, and $X_{n+1}$ ind. $X_1,...,X_n$)}\\
&= \mathbb{E}[\mathbb{E}[Y_n|X_1,...,X_{n-1}]|X_1,...,X_{n-1}] + X_n Z_{n-1} + 0
\quad\text{(induction hypothesis)}\\
&= \mathbb{E}[Y_n|X_1,...,X_{n-1}]+ X_n Z_{n-1}\\
&= Y_{n-1} + X_n Z_{n-1}\\
&= Y_n
\end{aligned}
\end{equation}
So
\[\forall n\geq 0,\mathbb{E}[Y_{n+1}|X_1,...,X_n]=Y_n\]
Then $Y_0,Y_1,...$ is a martingale with respect to $X_0,X_1,...$.

\clearpage
\section*{Question 5}
\paragraph{(a)}
Clearly, $M_n$ is a function of $X_1,...,X_n$.
Since $q/p>1,M_n>0,|M_n|=M_n$. Note that 
\[\mathbb{E}[X_i] = p-q\]
\begin{equation}
\begin{aligned}  
\mathbb{E}[|M_n|]&=\mathbb{E}\left[\left(\frac{q}{p}\right)^{X_1+...+X_n}\right]\\
&=\mathbb{E}\left[\prod_{i=1}^n \left(\frac{q}{p}\right)^{X_i}\right]\\
&=\prod_{i=1}^n \mathbb{E}\left[ \left(\frac{q}{p}\right)^{X_i}\right] \quad\text{(by independence)}\\
&< \infty
\end{aligned}
\end{equation}

\begin{equation}
\begin{aligned}  
\mathbb{E}[M_{n+1}|X_1,...,X_n]
&=\mathbb{E}\left[\left(\frac{q}{p}\right)^{X_1+...+X_{n+1}}|X_1,...,X_n\right]\\
&=\mathbb{E}\left[\left(\frac{q}{p}\right)^{X_1+...+X_n}
\left(\frac{q}{p}\right)^{X_{n+1}}|X_1,...,X_n\right]\\
&=\mathbb{E}\left[\left(\frac{q}{p}\right)^{X_{n+1}}\right]
\mathbb{E}\left[\left(\frac{q}{p}\right)^{X_1+...+X_n}|X_1,...,X_n\right]
\quad\text{($X_{n+1}$ ind. $X_1,...,X_n$)}\\
&=\left[p\left(\frac{q}{p}\right) + q\left(\frac{q}{p}\right)\right]
\left(\frac{q}{p}\right)^{X_1+...+X_n}\\
&=M_n
\end{aligned}
\end{equation}

So $M_i$ is a martingale with respect to the sequence $(X_i)$.

\clearpage
\paragraph{(b)}
We have $\mathbb{E}[T] < \infty$.
\begin{equation}
\begin{aligned}
\mathbb{E}[|M_{\min\{i+1,T\}} - M_{\min\{i,T\}}||X_1,...,X_i]
&= \mathbb{E}\left[\left|\left(\frac{q}{p}\right)^{X_1+...+X_{i+1}} 
- \left(\frac{q}{p}\right)^{X_1+...+X_i}\right||X_1,...,X_i\right]\\
&= \mathbb{E}\left[\left|\left(\frac{q}{p}\right)^{X_1+...+X_i}
\left(1-\left(\frac{q}{p}\right)^{X_{i+1}}\right)\right||X_1,...,X_i\right]\\
&\leq\left(\frac{q}{p}\right)^i \frac{q-p}{p}
\end{aligned}
\end{equation}
By the Optional Stopping Theorem, $\mathbb{E}[M_T]=\mathbb{E}[M_1] = 1$.

Let $a$ be the probability the gamber wins. Then
\[\mathbb{E}[M_T]=\mathbb{E}[(q/p)^{S_i}]
=a\left(\frac{q}{p}\right)^n + (1-a)\left(\frac{q}{p}\right)^{-k} = 1\]
We have
\[a=\frac{1-\left(\frac{p}{q}\right)^k}{\left(\frac{q}{p}\right)^n - \left(\frac{p}{q}\right)^k}\]

\paragraph{(c)}
Clearly, $Y_n$ is a function of $X_1,...,X_n$.
\[\mathbb{E}[|Y_n|] \leq i + (q-p)i < \infty\]
Note that 
\begin{equation}
\begin{aligned}
Y_{i+1} &= S_{i+1} + (q-p)(i+1)\\
&= S_i + X_{i+1} + (q-p)(i+1)\\
&= Y_i + X_{i+1} + q-p
\end{aligned}
\end{equation}

\begin{equation}
\begin{aligned}
\mathbb{E}[Y_{n+1}| X_1,...,X_n] &= \mathbb{E}[Y_n + X_{i+1} +q-p| X_1,...,X_n]\\
&= Y_n + \mathbb{E}[X_{i+1}] + q-p \quad\text{(independence)}\\
&= Y_n + p-q + q-p\\
&= Y_n
\end{aligned}
\end{equation}
Thus, $Y_i$ is a martingale with respect to the sequence $(X_i)$.

\clearpage
\paragraph{(d)}
We have $\mathbb{E}[T] < \infty$.
\begin{equation}
\begin{aligned}
\mathbb{E}[|Y_{\min\{i+1,T\}} - Y_{\min\{i,T\}}||X_1,...,X_i]
&= \mathbb{E}[|S_{i+1} + (q-p)(i+1)- S_i + (q-p)i||X_1,...,X_i]\\
&= \mathbb{E}[|X_{i+1} + q-p||X_1,...,X_i]\\
&\leq 1+q-p
\end{aligned}
\end{equation}
By the Optional Stopping Theorem, $\mathbb{E}[Y_T]=\mathbb{E}[Y_1] = 0$.

\[\mathbb{E}[Y_T]=\mathbb{E}[S_T+(q-p)T]=0\]
By linearity of expectation,
\[\mathbb{E}[S_T]+(q-p)\mathbb{E}[T]=0\]
\[\mathbb{E}[T]=\frac{\mathbb{E}[S_T]}{p-q}\]
Using part (b),
\begin{equation}
\begin{aligned}
\mathbb{E}[T]&=\frac{\mathbb{E}[S_T]}{p-q}\\
&=\frac{an+(1-a)(-k)}{p-q}\\
&=\frac{n+k-n\left(\frac{p}{q}\right)^k - k\left(\frac{q}{p}\right)^n}
{(p-q)\left[\left(\frac{q}{p}\right)^n-\left(\frac{p}{q}\right)^k\right]}
\end{aligned}
\end{equation}

\clearpage
\section*{Question 6}
\paragraph{(a)}
Let $X_i$ be the indicator random variable that an edge $e_i=(u,v)$ is monochromatic.
The probability that the edge is monochromatic is the probability that
both vertices are coloured the same so $\Pr(X_i=1)=1/2$.
Let $X(G)=\sum_i X_i$ be the total number of monochromatic edges in graph $G$.
Note that the number of edges is $|E|=m=dn/2$.
By linearity of expectation,
\[\mathbb{E}[X(G)]=\frac{dn}{2}\cdot\frac{1}{2}=\frac{dn}{4}\]

\paragraph{(b)}
Let $G=(V,E)$ and let $G_i=(V,E_i)$ where $E_i=\{e_1,e_2,...,e_i\}\subseteq E$.
Consider the Doob \textit{edge exposure} martingale 
\[Z_i = \mathbb{E}[X(G)|G_1,...,G_i]\]
Note that $Z_0=\mathbb{E}[X(G)]$.
By exposing a new edge, we cannot change the expected nuumber of monochromatic edges by more than 1 so
\[|Z_{i+1}-Z_i|\leq 1\]
We can apply the Azuma-Hoeffding inequality to get 
\[\Pr(|Z_m-Z_0|\geq \lambda \sqrt{m})\leq 2e^{-\lambda^2/2}\]
With high probability, the number of monochromatic edges is within $O(\sqrt{dn\log n})$ from its expected value.



\end{document}
